\documentclass[a4paper]{article}
\pdfoutput=1
\usepackage[utf8]{inputenc}
\usepackage{jheppub}
\usepackage{overpic}
\usepackage{subcaption}
\usepackage{fancybox}
\usepackage{float}
\usepackage{amsfonts}
\usepackage{amsmath}
\usepackage{amsbsy}
\usepackage{graphicx}
\usepackage{amssymb}
\usepackage{amsthm}
\usepackage{bm}
\usepackage{comment}
\usepackage{epsfig}
\usepackage{latexsym}
\usepackage{pdflscape}
\usepackage{color}
\usepackage{slashed}
\usepackage{tikz}
\usetikzlibrary{arrows.meta,decorations.markings}

\tikzset{
  mid arrow/.style={
    postaction={
      decorate,
      decoration={
        markings,
        mark=at position 0.5 with {\arrow{Stealth[length=6pt,width=6pt]}}
      }
    }
  }
}

\usepackage{cancel}
\usepackage{xcolor}
\usepackage{braket}
\usepackage{empheq}
\usepackage{graphicx} % Allows including images
\usepackage{tikz-3dplot}
% réglage de la vue 3D : élévation=20°, azimut=110°
\tdplotsetmaincoords{20}{110}
\usetikzlibrary{decorations.pathreplacing,calc}
\usetikzlibrary{decorations.pathmorphing}

\definecolor{blue5}{rgb}{0, 0, 0.66666}
\definecolor{red4}{rgb}{0.8, 0, 0}
\definecolor{navyblue}{rgb}{0, 0.2745, 0.67843}
\definecolor{cornflowerblue}{rgb}{0.388235, 0.6941176, 0.898039}
\definecolor{forestgreen}{rgb}{0.13, 0.55, 0.13}
\usetikzlibrary{decorations.text,decorations.pathmorphing}


\begin{document}



%%%%%%%%%%%%%%%%%%%%%%%%%%%%%%%%%%%%%%%%%%%%%%%%%%%%%%%%%%


% \begin{figure}[t!]
% \centering

% \begin{subfigure}{0.48\textwidth}
% \centering
% \resizebox{!}{2.8cm}{
% \begin{tikzpicture}[xscale=0.35, yscale=0.4]
% \hspace{0.3cm}

% \path[fill=cyan!10, opacity=0.25]
%   (-7.1,0)
% arc[start angle=180,end angle=270,x radius=1,y radius=3] 
%   .. controls (-5.,-2.2) and (-3.,-0.8) .. (-0.5,-0.8)
%   .. controls (0.1,-0.8) and (1.4,-2) .. (2.,-2)     
% arc[start angle=-90,end angle=-270,x radius=0.6,y radius=2]
%   .. controls (1.4,2) and (0.1,0.8) .. (-0.5,0.8)
%   .. controls (-3.,0.8) and (-5.,2.2) .. (-6.1,3)
% arc[start angle=90,end angle=180,x radius=1,y radius=3]
%   -- cycle;

% \path[fill=cyan!10, opacity=0.25]
%   (10.1,-3)
%   .. controls (9.,-2.2) and (7.,-0.8) .. (4.5,-0.8)
%   .. controls (3.9,-0.8) and (2.6,-2.) .. (2,-2)
% arc[start angle=-90,end angle=-270,x radius=0.6,y radius=2]
%   .. controls (2.6,2.) and (3.9,0.8) .. (4.5,0.8)
%   .. controls (7.,0.8) and (9.,2.2) .. (10.1,3)
% arc[start angle=90,end angle=-90,x radius=1,y radius=3]
%   -- cycle;

% \draw[thick, cornflowerblue, dotted] (-6.1,0) ellipse (1 and 3);
% \draw[thick,  cornflowerblue, dotted] (2,0) ellipse (0.6 and 2);
% \draw[thick,  cornflowerblue, dotted] (10.1,0) ellipse (1 and 3);


% \draw[thick, navyblue, decorate,
%       decoration={
%           text along path,
%           text={|\scriptsize|EAdS boundary},
%           text align=center,
%           text color=navyblue,
%           reverse path,
%           raise=3.5pt,
%           text format delimiters={|}{|}
%       }
% ]
% (-6.1,3) arc[start angle=90,end angle=270,x radius=1,y radius=3];


% \draw[thick, navyblue, decorate,
%       decoration={
%           text along path,
%           text={|\scriptsize|EAdS boundary},
%           text align=center,
%           text color=navyblue,
%           reverse path,
%           raise=0.pt,
%           text format delimiters={|}{|}
%       }
% ]
% (10.8,3) arc[start angle=90,end angle=-90,x radius=1,y radius=3];


% \draw[thick, cornflowerblue]
%   (-6,3) .. controls (-5.,2.2) and (-3,0.8) .. (-0.5,0.8)
%         .. controls (.1,0.8) and (1.4,2.) .. (2,2);

% \draw[thick, cornflowerblue]
%   (10,3) .. controls (9,2.2) and (7.,0.8) .. (4.5,0.8)
%         .. controls (3.9,0.8) and (2.6,2.) .. (2,2);

% \draw[thick, cornflowerblue]
%   (-6,-3) .. controls (-5.,-2) and (-3,-0.8) .. (-0.5,-0.8)
%          .. controls (0.1,-0.8) and (1.4,-2.) ..(2,-2);

% \draw[thick, cornflowerblue]
%   (10,-3) .. controls (9.,-2.2) and (7.,-0.8) .. (4.5,-0.8)
%         .. controls (3.9,-0.8) and (2.6,-2.) .. (2,-2);

% \draw[->,  line width=0.9pt,  opacity=0.2] (-6,0) -- (10.2,0);
% \node[above] at (9.7,0.) {\scriptsize\color{gray}$\tau$};

% \end{tikzpicture}
% }
% \end{subfigure}
% \hfill
% \begin{subfigure}{0.4\textwidth}
% \centering
% \resizebox{!}{2.8cm}{
% \begin{tikzpicture}[xscale=0.35, yscale=0.4]

% \path[fill=cyan!10, opacity=0.25]
%   (-7.1,0)
%   arc[start angle=180,end angle=270,x radius=1,y radius=3] 
%   .. controls (-5,-2.2) and (-3,-0.8) .. (-0.5,-0.8)
%   .. controls (0.1,-0.8) and (1.4,-2) .. (2,-2)     
%   arc[start angle=-90,end angle=-270,x radius=0.6,y radius=2]
%   .. controls (1.4,1.95) and (0.6,1.4) .. (0,0.8)
%   .. controls (-3,1) and (-4.8,2) .. (-6.1,3)
%   arc[start angle=90,end angle=180,x radius=1,y radius=3]
%   -- cycle;

% \path[fill=red, opacity=0.08]
%   (1.4,0)
%   arc[start angle=180,end angle=-90,x radius=0.6,y radius=2] 
%   .. controls (3.5,-2) .. (6.1,-3)                    
%   arc[start angle=-90,end angle=90,x radius=1,y radius=3]  
%   .. controls (3.5,2) .. (2,2)                        
%   arc[start angle=90,end angle=180,x radius=0.6,y radius=2]  
%   -- cycle;

% \pgfdeclareradialshading{inflationellipse}{\pgfpoint{0cm}{0cm}}{
%   color(0cm)=(red4!80);
%   color(0.6cm)=(orange!40);
%   color(0.9cm)=(yellow!10)
% }

% \shade[shading=inflationellipse, opacity=0.6] (2,0) ellipse (0.6 and 2);


% \draw[thick, navyblue, decorate,
%       decoration={
%           text along path,
%           text={|\scriptsize|EAdS boundary},
%           text align=center,
%           text color=navyblue,
%           reverse path,
%           raise=3.5pt,
%           text format delimiters={|}{|}
%       }
% ]
% (-6.1,3) arc[start angle=90,end angle=270,x radius=1,y radius=3];

% \draw[thick, red4, decorate,
%       decoration={
%           text along path,
%           text={|\scriptsize|Inflation},
%           text align=center,
%           text color=red4,
%           reverse path,
%           raise=0.pt,
%           text format delimiters={|}{|}
%       }
% ]
% (2.2,3) arc[start angle=95,end angle=-90,x radius=1,y radius=3];



% \draw[thick, red4, decorate,
%       decoration={
%           text along path,
%           text={|\tiny|Lorentzian evolution},
%           text align=center,
%           text color=red4,
%           reverse path,
%           raise=0.pt,
%           text format delimiters={|}{|}
%       }
% ]
% (6.8,3) arc[start angle=90,end angle=-90,x radius=1,y radius=3];



% \draw[thick, cornflowerblue, dotted] (-6.1,0) ellipse (1 and 3);
% \draw[thick, red4] (6.1,0) ellipse (1 and 3);
% \draw[thick, red4, dotted] (2,0) ellipse (0.6 and 2);

% \draw[thick, cornflowerblue]
%   (-6,3) .. controls (-5.,2.2) and (-3,0.8) .. (-0.5,0.8)
%         .. controls (.1,0.8) and (1.4,2.) .. (2,2);

% \draw[thick, red4] (2,2) .. controls (3.5,2) .. (6,3);

% \draw[thick, cornflowerblue]
%   (-6,-3) .. controls (-5.,-2) and (-3,-0.8) .. (-0.5,-0.8)
%          .. controls (0.1,-0.8) and (1.4,-2.) .. (2,-2);

% \draw[thick, red4] (2,-2) .. controls (3.5,-2) .. (6,-3);

% \draw[->,  line width=0.9pt,  opacity=0.2] (-6,0) -- (2,0);
% \node[above] at (1.1,0.) {\scriptsize\color{gray}$\tau$};

% \draw[->,  line width=0.9pt,  opacity=0.2] (2.2,0) -- (6.6,0);
% \node[above] at (6.,0.) {\scriptsize\color{gray}$t$};

% \end{tikzpicture}
% }
% \end{subfigure}
% \label{fig:wormholes}
% \end{figure}

% \begin{figure}[t]
% \centering

% % ============================================================
% % LEFT PANEL: complex t plane
% % ============================================================
% \begin{subfigure}[c]{0.42\textwidth}
% \centering

% \begin{tikzpicture}[
%     scale=0.68,
%     baseline=(current bounding box.center),
%     mid arrow/.style={
%         postaction={
%             decorate,
%             decoration={
%                 markings,
%                 mark=at position 0.5 with
%                 {\arrow{Stealth[length=6pt,width=6pt]}}
%             }
%         }
%     }
% ]

% \pgfdeclareradialshading{inflationellipse}{\pgfpoint{0cm}{0cm}}{
%   color(0cm)=(red4!80);
%   color(0.6cm)=(orange!40);
%   color(0.9cm)=(yellow!10)
% }

% % Axes
% \draw[->] (-3.1,0) -- (3.5,0)
% node[right] {$\Re\,(t)$};

% \draw[->] (0,-3.2) -- (0,3.2)
% node[above] {$\Im\,(t)$};

% % Upper Euclidean contour: (sigma,eta)=(+,+)
% \draw[
%     thick,
%     cornflowerblue,
%     postaction={
%         decorate,
%         decoration={
%             markings,
%             mark=at position 0.55 with
%             {\arrow{Stealth[length=6pt,width=6pt]}}
%         }
%     }
% ]
% (0,2.6) -- (0,0);

% % Lower alternative Euclidean contour: (sigma,eta)=(-,-)
% \draw[
%     thick,
%     cornflowerblue,
%     dashed,
%     postaction={
%         decorate,
%         decoration={
%             markings,
%             mark=at position 0.55 with
%             {\arrow{Stealth[length=6pt,width=6pt]}}
%         }
%     }
% ]
% (0,-2.6) -- (0,0);

% %Lorentzian segment
% \draw[
%     thick,
%     red4,
%     mid arrow
% ]
% (0,0) -- (2.8,0);


% % End points
% \fill[cornflowerblue] (0,2.6) circle (2.2pt);
% \fill[cornflowerblue] (0,-2.6) circle (2.2pt);
% \fill[red4] (2.8,0) circle (2.2pt);

% % Endpoint labels
% \node[left=5pt] at (0,2.6)
% {\small $\displaystyle \frac{i\pi}{2H}$};

% \node[left=5pt] at (0,-2.6)
% {\small $\displaystyle -\frac{i\pi}{2H}$};

% \node[below=5pt] at (2.8,0)
% {\small $t_f>0$};

% % Origin
% \node[below left=4pt] at (0,0)
% {\small $0$};

% % Region labels
% \node[cornflowerblue,left] at (-0.18,1.25)
% {\small Euclidean};

% \node[cornflowerblue,left,align=right] at (-0.18,-1.30)
% {\small alternative\\[-1mm]\small Euclidean};

% \node[red4,above] at (2.2,0.18)
% {\small Lorentzian};

% % Saddle labels
% \node[cornflowerblue,right] at (0.18,2.15)
% {\scriptsize $(\sigma,\eta)=(+,+)$};

% \node[cornflowerblue,right] at (0.18,-2.15)
% {\scriptsize $(\sigma,\eta)=(-,-)$};

% \end{tikzpicture}

% \end{subfigure}
% \hfill
% % ============================================================
% % RIGHT PANEL
% % ============================================================
% \begin{subfigure}[c]{0.54\textwidth}
% \centering

% \raisebox{-1.1\height}{
% \resizebox{0.6\linewidth}{!}{
% \begin{tikzpicture}[xscale=0.35,yscale=0.4]

% % Euclidean hemisphere
% \path[fill=cyan!10, opacity=0.3]
%     (2,2)
%     arc[  start angle=90,  end angle=270,  x radius=2,  y radius=2 ]
%     -- (2,-2) -- cycle;

% \draw[thick,cornflowerblue]
%     (2,2)
%     arc[ start angle=90, end angle=270,  x radius=2,  y radius=2 ];

% % no-boundary label
% \draw[ thick, navyblue, decorate, decoration={
%         text along path,
%         text={|\tiny|no-boundary},
%         text align=center,
%         text color=navyblue,
%         reverse path,
%         raise=3.5pt,
%         text format delimiters={|}{|}
%     }
% ]
% (2,2)
% arc[start angle=90, end angle=270, x radius=2, y radius=2];

% % Lorentzian region
% \path[fill=red,opacity=0.08]
%  (1.4,0)
%  arc[ start angle=180, end angle=-90, x radius=0.6, y radius=2]
%  .. controls (3.5,-2) .. (6.1,-3)
%  arc[ start angle=-90, end angle=90, x radius=1, y radius=3
%  ]
%  .. controls (3.5,2) .. (2,2)
%  arc[ start angle=90,  end angle=180,  x radius=0.6,  y radius=2 ]
%  -- cycle;

% % Matching surface
% \shade[
%     shading=inflationellipse,
%     opacity=0.6
% ]
% (2,0) ellipse (0.6 and 2);

% % Lorentzian boundaries
% \draw[thick,red4]
% (6.1,0) ellipse (1 and 3);

% \draw[thick,red4,dotted]
% (2,0) ellipse (0.6 and 2);

% \draw[thick,red4]
% (2,2)
% .. controls (3.5,2) ..
% (6,3);

% \draw[thick,red4]
% (2,-2)
% .. controls (3.5,-2) ..
% (6,-3);

% % Inflation label
% \draw[thick, red4, decorate, decoration={
%         text along path,
%         text={|\scriptsize|Inflation},
%         text align=center,
%         text color=red4,
%         reverse path,
%         raise=0pt,
%         text format delimiters={|}{|}
%     }
% ]
% (2.2,3)
% arc[
%     start angle=95,
%     end angle=-90,
%     x radius=1,
%     y radius=3
% ];

% % Lorentzian evolution label
% \draw[ thick, red4, decorate, decoration={
%         text along path,
%         text={|\tiny|Lorentzian evolution},
%         text align=center,
%         text color=red4,
%         reverse path,
%         raise=0pt,
%         text format delimiters={|}{|}
%     }
% ]
% (6.8,3)
% arc[ start angle=90, end angle=-90,  x radius=1,  y radius=3
% ];

% % Horizontal Euclidean time line
% \draw[->, line width=0.9pt, opacity=0.2]
% (0,0) -- (2,0);

% \node[above] at (1.1,0)
% {\scriptsize\color{gray}$\tau$};

% % Horizontal Lorentzian time line
% \draw[
%     ->,
%     line width=0.9pt,
%     opacity=0.2
% ]
% (2.2,0) -- (6.6,0);

% \node[above] at (6,0)
% {\scriptsize\color{gray}$t$};

% \end{tikzpicture}
% }}
% \end{subfigure}

% \label{fig:no-boundary-contour}

% \end{figure}

% \begin{figure}[htb]
% \centering
% \begin{tikzpicture}


% \path[fill=cyan!10, opacity=0.25]
%   (-7.1,0)
%   arc[start angle=180,end angle=270,x radius=1,y radius=3] 
%   .. controls (-4.8,-2) and (-3,-1) .. (0,-0.8)
%   .. controls (0.6,-1.4) and (1.4,-1.95) .. (2,-2)     
%   arc[start angle=-90,end angle=-270,x radius=0.6,y radius=2]
%    % change this angle=270 to  angle=90
%   .. controls (1.4,1.95) and (0.6,1.4) .. (0,0.8)
%   .. controls (-3,1) and (-4.8,2) .. (-6.1,3)
%   arc[start angle=90,end angle=180,x radius=1,y radius=3]
%   -- cycle;

% \path[fill=red, opacity=0.08]
%   (1.4,0)
%   arc[start angle=180,end angle=-90,x radius=0.6,y radius=2] 
%    % change this angle=-90 to  angle=270
%   .. controls (3.5,-2) .. (6.1,-3)                    
%   arc[start angle=-90,end angle=90,x radius=1,y radius=3]  
%   .. controls (3.5,2) .. (2,2)                        
%   arc[start angle=90,end angle=180,x radius=0.6,y radius=2]  
%   -- cycle;

% % \shade[inner color=red4!60,  outer color=orange!40, opacity=0.6] 
% %   (2,0) ellipse (0.6 and 2);

% \pgfdeclareradialshading{inflationellipse}{\pgfpoint{0cm}{0cm}}{
%   color(0cm)=(red4!80);
%   color(0.6cm)=(orange!40);
%   color(0.9cm)=(yellow!10)
% }
% \shade[shading=inflationellipse, opacity=0.6] (2,0) ellipse (0.6 and 2);



% % horizontal axis
% \draw[->] (-6,0) -- (6.2,0);

% % Points
% \filldraw[fill=navyblue] (-6,0) circle (2pt) node[below] {$\begin{array}{c} \color{navyblue}{z_n^E \to 0}\\ \tau \to -\infty\end{array}$};
% \filldraw[fill=red4] (2,0) circle (2pt);
% \filldraw[fill=red4] (6,0) circle (2pt) node[below] {$\begin{array}{c} \color{red4}{z_n^L \to 0} \\ t \to +\infty \end{array}$};


% \draw[thick, red4, decorate, 
%       decoration={text along path, text={Lorentzian evolution}, text align=center, text color=red4}] 
%       (2,2.1) .. controls  (3.5,2.1) .. (6,3.1);
% \draw[thick, navyblue, decorate, 
%       decoration={text along path, 
%                   text={Euclidean past}, 
%                   text align=center,
%                   text color=navyblue}] 
%       (-5.9,3.1) .. controls (-4.8,2.1) and (-3,1.1) .. (0,0.8);
% \draw[thick, navyblue, decorate,
%       decoration={
%       text along path,
%                   text={EAdS boundary},
%                   text align=center,
%                   text color=navyblue, reverse path, raise=3.5pt}
%                 ] 
%        (-6,3) .. controls (-7.4,2.8) and (-7.4,-2.8)  .. (-6,-3);
% \draw[thick, red4, decorate,
%       decoration={
%       text along path,
%                   text={Inflation},
%                   text align=center,
%                   text color=red4
%                    ,  raise=3pt
%                   }
%                 ] 
%        (2,2) .. controls (2.8,1.8) and (2.8,-1.8)  .. (2,-2);


% % Regions
% \node[above] at (-4,0.5) {\Large\color{cornflowerblue}{I}};
% \node[above] at (0.8,0.5) {\Large\color{navyblue}{II}};


% % Legend points
% \node[below] at (2.2,-2.2) {$t =0 = \tau$};

% % Waves

% \draw[->, decorate, decoration={snake, amplitude=0.4mm, segment length=3mm}, thick, navyblue] (-5,-1) -- (-4,-1) node[midway, above] {$\Phi^\star_{n}(\tau)$};
% \draw[->, decorate, decoration={snake, amplitude=0.4mm, segment length=3mm}, thick, navyblue] (-5,-1.3) -- (-4,-1.3);


% \draw[->, decorate, decoration={snake, amplitude=0.4mm, segment length=3mm}, thick, red4] (0.5,-1) -- (1.5,-1) node[midway, above] {$\Phi^\star_{n}(\tau)$};
% \draw[->, decorate, decoration={snake, amplitude=0.4mm, segment length=3mm}, thick, red4] (0.5,-1.3) -- (1.5,-1.3);


% \draw[->, decorate, decoration={snake, amplitude=0.4mm, segment length=3mm}, thick, red4] (3.7,-1) -- (4.7,-1) node[midway, above] {$\Phi^\star_{n}(t)$};
% \draw[->, decorate, decoration={snake, amplitude=0.4mm, segment length=3mm}, thick, red4] (3.7,-1.3) -- (4.7,-1.3);

% \draw[thick, cornflowerblue, dotted] (-6.1,0) ellipse (1 and 3);   %
% \draw[thick, red4] (6.1,0) ellipse (1 and 3);   % 
% \draw[thick, red4, dotted] (2,0) ellipse (0.6 and 2);

% % 
% % \draw[thick, navyblue]
% %   (-6,3) .. controls (-4.8,2) and (-3,1) .. (0,0.8)
% %         .. controls   (0.6,1.4) and (1.4,1.95) .. (2,2);

% % \draw[thick, navyblue]
% %   (-6,-3) .. controls (-4.8,-2) and (-3,-1) .. (0,-0.8)
% %          .. controls (0.6,-1.4) and (1.4,-1.95) ..(2,-2);

% \draw[thick, cornflowerblue]
%   (-6,3) .. controls (-4.8,2) and (-3,1) .. (0,0.8);

% \draw[thick, cornflowerblue]
%   (-6,-3) .. controls (-4.8,-2) and (-3,-1) .. (0,-0.8);

% % Δεξί (γαλάζιο) κομμάτι
% \draw[thick, dashed, navyblue]
%   (0,0.8) .. controls (0.6,1.4) and (1.4,1.95) .. (2,2);

% \draw[thick, dashed, navyblue]
%   (0,-0.8) .. controls (0.6,-1.4) and (1.4,-1.95) .. (2,-2);


% \draw[thick, red4] (2,2) .. controls  (3.5,2) .. (6,3);
% \draw[thick, red4] (2,-2) .. controls (3.5,-2) .. (6,-3);

% \draw[<->,thick, navyblue] (-6,-3.8) -- (-4,-3.8);
% \node[below] at (-5,-4) {\color{navyblue} \text{Bessel}};

% \draw[<->,thick, navyblue] (-4.5,-3.4) -- (2,-3.4);
% \node[below] at (-1,-3.6) {\color{navyblue} \text{WKB approximation}};

% \draw[<->,thick, red4] (2,-3.4) -- (5,-3.4);
% \node[below] at (3.3,-3.4) {$\begin{array}{c} {\color{red4}\text{WKB}}\\{\color{red4} \text{Approximation}}\end{array}$};

% \draw[<->,thick, red4] (4.5,-3.8) -- (6,-3.8);
% \node[below] at (5.3,-4) {\color{red4} \text{Bessel}};

% \draw[-, dashed, green] (2, -2.8) -- (2, -4.2) ;


% \end{tikzpicture}
% \caption{Schematic view of the WKB matching problem.}
% \label{fig:waves_WKB}
% \end{figure}

%   \begin{equation*}  
% U(P|\sigma) \sim \begin{tikzpicture}[baseline={([yshift=-.5ex]current bounding box.center)}, scale=0.7]
% \draw[fill=cyan!40!white,opacity=0.5] (0,0) ellipse (1.5 and 1.5);
% \draw[fill] (0,0) circle (0.06);
% \node at (1.9,0.5) {\small $\sigma$};
% \node at (0.5,0) {\small $\mathcal{V}_{P}$};
% \end{tikzpicture} \qquad 
% R(P,\tilde{P}|\sigma) \sim \begin{tikzpicture}[baseline={([yshift=-.5ex]current bounding box.center)}, scale=0.7]
% \draw[fill=cyan!40!white,opacity=0.5] (0,0) ellipse (1.5 and 1.5);
% \draw[fill] (1.5,0) circle (0.06); 
% \node at (0,-1.985) {};
% \node at (0,1.8) {\small $\sigma$};
% \draw[fill] (0,0) circle (0.06);
% \node at (-0.5,0) {\small $\mathcal{V}_P$};
% \node at (2.1,0) {\small $\Psi^{\sigma \sigma}_{\tilde{P}}$};
% \end{tikzpicture}
% \vspace{-0.4cm}
% \end{equation*}

% \begin{equation*}
% C^{\sigma_1 \sigma_2 \sigma_3}_{\tilde{P}_1 \tilde{P}_2 \tilde{P}_3} \sim  \begin{tikzpicture}[baseline={([yshift=-.5ex]current bounding box.center)}, scale=0.7]
% \draw[fill=cyan!40!white,opacity=0.5] (0,0) ellipse (1.5 and 1.5);
% \draw[fill] (0,1.5) circle (0.06);
% \draw[fill] (-1.3,-0.75) circle (0.06); 
% \draw[fill] (1.3,-0.75) circle (0.06); 
% %\draw (-1.3,-0.75) -- (0,0) -- (1.3,-0.75);
% %\draw (0,0) -- (0,1.5); 
% \node at (-2,-0.8) {\small $\Psi^{\sigma_3 \sigma_1}_{\tilde{P}_{3}}$};
% \node at (2.1,-0.8) {\small $\Psi^{\sigma_2 \sigma_3}_{\tilde{P}_{2}}$};
% \node at (0,-1.85) {\small $\sigma_3$};
% \node at (-1.9,1) {\small $\sigma_1$};
% \node at (1.9,1) {\small $\sigma_2$};
% \node at (0,1.9) {\small $\Psi^{\sigma_1 \sigma_2}_{\tilde{P}_{1}}$};
% \end{tikzpicture}
% \vspace{-0.4cm} 
% \end{equation*}

\end{document}

